\begin{figure*}[t]
\centering
\resizebox{\textwidth}{!}{% Fig. 2: architecture of the method (schematic in TikZ, v4 2026-10-23). Spec: docs/writing/v2_figure_specs.md (Fig. 2).
% Semantic colors: obscol (observation), predcol (prediction / world), actcol (action), gtcol (neutral).
% The picture is about 17.4 cm wide and is scaled to \textwidth, so \scriptsize prints at about 7 pt.
\begin{tikzpicture}[
  yscale=0.92,
  font=\sffamily\scriptsize,
  >={Stealth[length=4pt,width=3pt]},
  box/.style={draw=gtcol!70, rounded corners=2pt, align=center, inner sep=2.5pt, fill=white, text width=2.85cm},
  lane/.style={rounded corners=2pt, align=left, inner sep=2.5pt, minimum height=0.60cm, text width=3.95cm},
  flowarrow/.style={->, thick, gtcol!80},
  readout/.style={->, thick, predcol, dashed},
]
% ======================= Titles =======================
\node[font=\sffamily\footnotesize\bfseries, anchor=west] at (0,5.12) {Inputs at capture time $\tau_t$};
\node[font=\sffamily\footnotesize\bfseries, anchor=west] at (4.55,5.12) {World--action transformer};
\node[font=\sffamily\footnotesize\bfseries, anchor=west] at (11.75,5.12) {World outputs (training only)};
% ======================= Inputs =======================
\draw[gtcol!70, fill=gtcol!8, rounded corners=1pt] (0.05,3.92) rectangle (1.10,4.80);
\fill[gtcol!25] (0.05,3.92) -- (0.42,4.30) -- (0.70,4.12) -- (1.10,4.50) -- (1.10,3.92) -- cycle;
\fill[obscol] (0.48,4.40) rectangle (0.66,4.56);
\node[box, anchor=west] (img) at (1.22,4.36) {Current RGB-D frame\\frozen \cfgVisualEncoder\\\cfgPatches{} patches + metric XYZ};
\draw[obscol!35, line width=0.6pt] (0.12,2.72) .. controls (0.45,3.37) and (0.75,2.77) .. (1.05,3.30);
\foreach \x/\y/\o in {0.12/2.72/30, 0.42/3.05/50, 0.74/2.97/75, 1.05/3.30/100}{
  \begin{scope}[shift={(\x,\y)}]
    \draw[obscol!\o, line width=0.7pt, ->] (0,0) -- (0.20,0.06);
    \draw[obscol!\o, line width=0.7pt, ->] (0,0) -- (-0.05,0.20);
  \end{scope}}
\node[anchor=north west, text=obscol, inner sep=1pt] at (-0.02,2.45) {$-$\cfgHistorySpan{} to 0};
\node[box, anchor=west, draw=obscol] (hist) at (1.22,3.02) {Target SE(3) history\\\cfgHistoryLen{} steps: pose +\\interval motion};
\node[box, anchor=west] (robot) at (1.22,1.86) {Robot state history\\\cfgHistoryLen{} steps};
% ======================= Transformer =======================
\draw[gtcol!60, rounded corners=4pt, fill=gtcol!4] (4.55,1.27) rectangle (11.35,4.88);
\draw[predcol!70, rounded corners=3pt, fill=predcol!5, dashed] (4.75,2.80) rectangle (9.18,4.84);
\node[anchor=south west, text=predcol!80!black, inner sep=1pt] at (4.80,2.82) {world stream};
\node[lane, fill=obscol!12, draw=obscol] (laneS) at (6.95,4.38) {\textbf{Scene} $\mathcal S$: \cfgPatches{} patch,\\\cfgHistoryLen{} target, \cfgHistoryLen{} robot tokens};
\node[lane, fill=predcol!15, draw=predcol] (laneM) at (6.95,3.48) {\textbf{Summary} $\mathcal M$: \cfgMetaTokens{}\\learnable tokens};
\node[lane, fill=actcol!12, draw=actcol] (laneA) at (6.95,1.90) {\textbf{Action} $\mathcal A$: \cfgChunk{} commands,\\conditional flow matching};
\node[align=center, text=gtcol, text width=2.05cm] at (10.35,2.45) {shared attention,\\stream FFNs\\($\times$\cfgLayers)};
\draw[->, thick, gtcol!70] (laneS.south) -- (laneM.north);
\draw[->, thick, actcol] (laneM.south) -- node[right, text=actcol!80!black, align=left] {reads summary\\tokens only} (laneA.north);
% input arrows
\draw[flowarrow] (img.east) -- (laneS.west |- img.east);
\draw[flowarrow] (hist.east) -- ++(0.18,0) |- ([yshift=-3pt]laneS.west);
\draw[flowarrow] (robot.east) -- ++(0.30,0) |- ([yshift=-8pt]laneS.west);
% ======================= Outputs =======================
\draw[predcol!60, line width=0.6pt] (11.85,3.92) .. controls (12.65,4.72) and (13.30,3.92) .. (14.20,4.57);
\foreach \x/\y/\o in {11.85/3.92/40, 12.40/4.30/55, 12.98/4.24/70, 13.60/4.25/85, 14.20/4.57/100}{
  \begin{scope}[shift={(\x,\y)}]
    \draw[predcol!\o!gtcol, line width=0.8pt, ->] (0,0) -- (0.24,0.05);
    \draw[predcol!\o!gtcol, line width=0.8pt, ->] (0,0) -- (-0.04,0.24);
  \end{scope}}
\node[anchor=west, align=left] at (14.50,4.26) {\textbf{Target SE(3)}\\$(\vect d_h,\boldsymbol\phi_h)$, $h\le\cfgHorizons$};
\foreach \i in {0,...,4}{\foreach \j in {0,...,2}{
  \pgfmathsetmacro{\dx}{0.18+0.05*\j}
  \draw[predcol!70, ->, line width=0.55pt] ({11.90+0.50*\i},{2.62+0.27*\j}) -- ++(\dx,{0.06*\i-0.08});}}
\node[anchor=west, align=left] at (14.50,3.05) {\textbf{Sparse 3D flow}\\$F_h$ + visibility, \cfgContextGrid\\\textcolor{gtcol}{\footnotesize + auxiliary targets}};
\draw[readout] (9.18,4.22) -- (11.72,4.22);
\draw[readout] (9.18,3.30) -- (11.24,3.30) |- (11.72,2.86);
\node[anchor=south, text=predcol!80!black] at (10.45,4.22) {readout heads};
% attention mask (rows read columns): S, M, A
\def\mx{12.30}\def\my{1.12}\def\cs{0.30}
\foreach \r/\rn in {0/S,1/M,2/A}{
  \node[anchor=east] at (\mx-0.03,{\my+(2-\r)*\cs+0.5*\cs}) {$\mathcal{\rn}$};
  \node[anchor=south] at ({\mx+\r*\cs+0.5*\cs},{\my+3*\cs-0.04}) {$\mathcal{\rn}$};}
\foreach \r/\c/\col in {0/0/obscol, 1/0/predcol, 1/1/predcol, 2/1/actcol, 2/2/actcol}{
  \fill[\col!45] ({\mx+\c*\cs},{\my+(2-\r)*\cs}) rectangle ++(\cs,\cs);}
\draw[gtcol!80, shift={(\mx,\my)}] (0,0) grid[step=\cs] ({3*\cs},{3*\cs});
\node[anchor=west, align=left] at ({\mx+3*\cs+0.18},{\my+1.5*\cs}) {attention mask $\mathcal G$\\(row reads column)};
% ======================= Physical time =======================
% 1.25 cm per 50 ms; capture at x = 5.00. The query ends at 80 ms, so the chunk takes
% over at the 100 ms boundary (x = 7.50) from command 2; commands 0-1 have expired.
\fill[actcol!14] (7.50,0.43) rectangle (10.00,0.67);
\draw[->, gtcol] (4.55,0.55) -- (17.35,0.55);
\node[anchor=north east, text=gtcol] at (17.35,0.47) {physical time};
\draw[gtcol, line width=0.8pt] (5.00,0.42) -- (5.00,0.70);
\node[anchor=north east] at (5.02,0.42) {$\tau_t$};
\foreach \h in {1,...,8}{\draw[predcol, line width=0.7pt] ({5.0+1.25*\h},0.48) -- ({5.0+1.25*\h},0.62);}
\node[anchor=north, text=actcol!75!black] at (8.75,0.40) {executed};
\node[anchor=north, text=predcol] at (13.30,0.40) {forecast horizons $\tau_t+\delta_h$, $\delta_h=h\cdot$\cfgHorizonStep};
\foreach \k in {0,1}{
  \draw[gtcol!70, dashed] ({4.86+1.25*\k},0.70) rectangle ++(0.28,0.26);}
\foreach \k in {2,...,7}{
  \pgfmathtruncatemacro{\shade}{\k<4 ? 85 : 30}
  \fill[actcol!\shade] ({4.86+1.25*\k},0.70) rectangle ++(0.28,0.26);}
\node[anchor=east, text=gtcol] at (4.78,0.83) {expired};
\draw[<->, actcol] (5.00,0.30) -- (7.50,0.30) node[midway, below=-1pt] {latency $\ell_t$};
\node[anchor=west, text=actcol] at (14.02,0.83) {$\tapp_{t,j}$};
\draw[flowarrow, actcol] (laneA.east) -- (9.60,1.90) -- (9.60,1.10) -| (7.50,0.99);
% replan loop
\draw[->, dashed, gtcol!80] (7.50,0.43) -- (7.50,-0.05) -- (0.25,-0.05) -- (0.25,1.52);
\node[anchor=south west, text=gtcol] at (0.40,-0.04) {execute $\to$ observe $\to$ replan};
\end{tikzpicture}
}
\caption{\textbf{\ours。}世界流包含 \cfgSceneTokens{} 个场景 token（冻结编码器给出的当前图像 patch，以及目标度量 SE(3) 位姿与机器人状态的短历史）和 \cfgMetaTokens{} 个摘要 token。训练时，轻量读出头把世界流读出为物理时域 $\tau_t+\delta_h$ 上按角色区分精度的 4D 预测：目标的 SE(3) 轨迹，以及覆盖整个场景、位于 \cfgContextGrid{} 网格上带可见性的稀疏 3D 流（虚线：仅训练）。动作流以流匹配生成 \cfgChunk{} 条指令组成的动作块。注意力掩码 $\mathcal G$ 只允许场景 token 读取场景、摘要 token 读取场景与摘要、动作 token 读取摘要与动作。策略在每次查询结束时或之后的首个控制边界重新观测并重新规划，新动作块从对应该边界的那条指令开始接管执行。}
\label{fig:architecture}
\end{figure*}